
\medskip
\phantomsection\label{lem:sga3-VIB-5.6.2.0}%
\noindent\textbf{Lemma 5.6.2.0.---}%
\footnote{\label{note:sga3-VIB-42}Editor's note: we have added this
lemma, communicated by O.\ Gabber. It makes it possible to simplify
the proof of \SGARefLink{cor:sga3-VIB-5.6.2}{5.6.2} and to prove \hyperref[thm:sga3-VIB-5.3]{Theorem 5.3}, as well as \hyperref[cor:sga3-VIB-5.4]{5.4}, in a
more general form; see \hyperref[thm:sga3-VIB-5.7]{5.7} and \hyperref[sec:sga3-VIB-5.8]{5.8} below.}
\emph{Let \(k\) be a field and let \(G\) be a nonempty
\(k\)-scheme of finite type.}
\begin{enumerate}
\item[\textup{\SGAInlineRefTarget{lem:sga3-VIB-5.6.2.0:item:i}(i)}] \emph{Let \(K\) be an extension of \(k\), and
  let \(W\) be a dense open subset of \(G_K\). Denote by \(\pi\)
  the projection \(G_K\to G\). Then}
  \[
    U=\bigl\{g\in G\mathrel{\big|}W\text{ contains every maximal
      point of }\pi^{-1}(g)\bigr\}
  \]
  \emph{is a dense open subset of \(G\). \textup{(N.B. If \(g\)
  is a closed point of \(G\) belonging to \(U\), then
  \(\pi^{-1}(g)\subset W\).}}

\item[\textup{\SGAInlineRefTarget{lem:sga3-VIB-5.6.2.0:item:ii}(ii)}] \emph{Suppose in addition that \(G\) is
  geometrically irreducible. Let}
  \[
    \mu:G\times X\longrightarrow Y
  \]
  \emph{be a morphism of \(k\)-schemes, let \(x\) be a point of
  \(X\), and let \(\Omega\) be an open subset of \(Y\) such that}
  \[
    \mu_x^{-1}\bigl(\Omega_{\kappa(x)}\bigr)\ne\varnothing,
  \]
  \emph{where \(\mu_x\) denotes the morphism}
  \[
    G_{\kappa(x)}\longrightarrow Y_{\kappa(x)},
    \qquad g\longmapsto\mu(g,x).
  \]
  \emph{Then}
  \[
    U=\bigl\{g\in G\mathrel{\big|}\mu\text{ sends every maximal
      point of }g\times x\text{ into }\Omega\bigr\}
  \]
  \emph{is a dense open subset of \(G\). For every closed point
  \(g\) of \(G\) belonging to \(U\), one has}
  \[
    \mu(g\times x)\subset\Omega,
  \]
  \emph{and hence}
  \[
    \mu(g',x)\in\Omega_{\kappa(x)}
    \quad\text{\textup{(respectively, }}\,
    \mu(g,x')\in\Omega_{\kappa(g)}\text{\textup{)}},
  \]
  \emph{for every point \(g'\in G_{\kappa(x)}\) above \(g\)
  \textup{(respectively, \(x'\in X_{\kappa(g)}\) above \(x\))}.}

\item[\textup{\SGAInlineRefTarget{lem:sga3-VIB-5.6.2.0:item:iii}(iii)}] \emph{Suppose that \(G=H^0\), where \(H\)
  is a \(k\)-group scheme locally of finite type, acting on a
  nonempty \(k\)-scheme \(X\) in such a way that the morphism}
  \[
    H\times_S X\longrightarrow X\times_S X,
    \qquad (h,x)\longmapsto(hx,x)
  \]
  \emph{is surjective. Let \(U\) be an open subset of \(X\).
  Then:}
  \begin{enumerate}
  \item[\textup{\SGAInlineRefTarget{lem:sga3-VIB-5.6.2.0:item:a}(a)}] \emph{\(G\cdot U\) is an open subset of
    \(X\), equal to the union of the irreducible components of
    \(X\) whose generic point belongs to \(U\).}
  \item[\textup{\SGAInlineRefTarget{lem:sga3-VIB-5.6.2.0:item:b}(b)}] \emph{Suppose that \(U\) is dense in every
    irreducible component of \(X\). Then, for every finite subset
    \(A\) of \(X\), there exists a closed point \(g\in G\) such
    that}
    \[
      g\cdot A'\subset U_{\kappa(g)},
    \]
    \emph{where \(A'\) is the inverse image of \(A\) in
    \(X_{\kappa(g)}\). In particular, the morphism}
    \[
      G\times U\longrightarrow X
    \]
    \emph{is surjective.}
  \end{enumerate}
\end{enumerate}

\emph{Proof.}
\textup{\SGARefLink{lem:sga3-VIB-5.6.2.0:item:i}{(i)}}
Let \(\bar{k}\) be a separable closure of \(k\), and let \(L\) be
an extension of \(k\) containing a copy of \(\bar{k}\) and of
\(K\). Denote by \(\pi_L\) (respectively, \(\phi\)) the projection
\[
  G_L\longrightarrow G
  \qquad\text{(respectively, \(G_L\longrightarrow G_K\)).}
\]
For every \(g\in G\), the map \(\phi\) sends the maximal points of
\(\pi_L^{-1}(g)\) surjectively onto those of \(\pi^{-1}(g)\).
Consequently,
\[
  U=\bigl\{g\in G\mathrel{\big|}W_L\text{ contains every maximal
    point of }\pi_L^{-1}(g)\bigr\}.
\]
Thus, replacing \(K\) by \(L\), we are reduced to the case in
which \(K\) contains \(\bar{k}\).

Since the projection
\[
  p:G_K\longrightarrow G_{\bar{k}}
\]
is surjective and open, \(V=p(W)\) is a dense open subset of
\(G_{\bar{k}}\). Moreover,
\[
  p^{-1}(g)=\operatorname{Spec}
    \bigl(\kappa(g)\otimes_{\bar{k}}K\bigr)
\]
is irreducible for every \(g\in G_{\bar{k}}\) (cf.\
EGA~IV\(_2\), 4.3.3). Hence, for every \(g\in V\), the generic
point of \(p^{-1}(g)\) belongs to \(W\).

On the other hand, set
\[
  \mathscr G=\operatorname{Gal}(\bar{k}/k).
\]
The projection \(q:G_{\bar{k}}\to G\) is surjective, and
\(\mathscr G\) acts transitively on the fibers of \(q\).
Therefore
\[
  U=q(V'),
\]
where \(V'\) is the intersection of the
\(\mathscr G\)-conjugates of \(V\).

Let \(Z\) be the closed subset \(G_{\bar{k}}-V\), endowed with
its reduced closed subscheme structure. Since \(G_{\bar{k}}\),
and hence \(Z\), is of finite presentation over \(\bar{k}\), the
scheme \(Z\) is obtained by base change from a reduced closed
subscheme
\[
  Z_1\subset G\otimes_k k_1
\]
for some finite Galois extension \(k_1\) of \(k\). Thus the
\(\mathscr G\)-conjugates of \(Z\) are finite in number, and their
union is again a nowhere-dense closed subset \(F\) of
\(G_{\bar{k}}\). It follows that
\[
  V'=G_{\bar{k}}-F
\]
is a dense open subset of \(G_{\bar{k}}\). Since
\(q:G_{\bar{k}}\to G\) is surjective and open, \(U=q(V')\) is
therefore a dense open subset of \(G\).

Moreover, for every closed point \(g\) of \(G\), the fiber
\(\pi^{-1}(g)\) consists of finitely many closed points of \(G_K\).
Thus, if \(g\in U\), then
\[
  \pi^{-1}(g)\subset W.
\]
This proves \textup{\SGARefLink{lem:sga3-VIB-5.6.2.0:item:i}{(i)}}, and \textup{\SGARefLink{lem:sga3-VIB-5.6.2.0:item:ii}{(ii)}} follows from it.

Assertion \textup{\SGARefLink{lem:sga3-VIB-5.6.2.0:item:iii}{(iii)}\SGARefLink{lem:sga3-VIB-5.6.2.0:item:a}{(a)}} was proved in VI\(_A\), \SGARefLink{prop:sga3-VIA-2.6.4}{2.6.4}.
Finally, if \(U\) is dense in every irreducible component of
\(X\), then
\[
  X=G\cdot U.
\]
Hence, for every \(x\in X\),
\[
  \mu_x^{-1}\bigl(U_{\kappa(x)}\bigr)
\]
is a nonempty open subset of \(G_{\kappa(x)}\), and
\textup{\SGARefLink{lem:sga3-VIB-5.6.2.0:item:iii}{(iii)}\SGARefLink{lem:sga3-VIB-5.6.2.0:item:b}{(b)}} follows from \textup{\SGARefLink{lem:sga3-VIB-5.6.2.0:item:ii}{(ii)}}.
